% =====================================================================
\chapter{Analyse de l'existant et recentrage}
\label{ch:existant}
% =====================================================================

Ce chapitre raconte la décision la plus importante du stage, et la moins
confortable : supprimer une grande partie de ce qui avait été construit.

\section{État de départ : un prototype fonctionnel}

Au début du stage, le dépôt contenait déjà une plateforme climatique pour
l'immobilier français. Elle interrogeait quatre sources publiques — Géorisques,
BDNB, Copernicus, Open-Meteo — et présentait les résultats dans un tableau de
bord. Techniquement, cela fonctionnait : les connecteurs répondaient, les
données s'affichaient, une couche de notation produisait un score, une
visualisation 3D extrudait les bâtiments.

La couche géométrique la plus difficile était même déjà écrite : un parseur GML
3.2 fait main et un test point-in-polygon par lancer de rayon, sans dépendance
à \texttt{shapely}.

\section{Le diagnostic : quatre symptômes, une cause}

En auditant le prototype, quatre problèmes sont apparus.

\subsection{La redondance avec les portails sources}

C'est le problème central. Pour chaque aléa, le prototype affichait ce que
Géorisques affiche déjà, à la même résolution, gratuitement. Pour chaque
bâtiment, il affichait ce que la BDNB affiche déjà. Le test appliqué à chaque
fonctionnalité a été formulé ainsi :

\begin{quote}
\itshape
L'acheteur pourrait-il obtenir ceci directement depuis la source ?
\end{quote}

Si la réponse est oui, la fonctionnalité ne justifie pas son existence. Un
agrégateur qui reproduit ses sources n'apporte pas de valeur : il apporte une
dépendance supplémentaire.

\subsection{Le score composite}

Le prototype produisait un score global par bâtiment, agrégé à partir des
aléas. C'est intuitivement satisfaisant — un chiffre, facile à comparer, facile
à afficher. Mais c'est précisément ce que le métier refuse (chapitre~\ref{ch:metier}) :

\begin{itemize}[leftmargin=1.4em]
  \item un score tiers engage la responsabilité de son producteur ;
  \item un score opaque est soit ignoré, soit reconstruit par l'actuaire ;
  \item un score mélange des aléas non comparables (inondation et sismicité ne
        se compensent pas) ;
  \item sous Solvabilité~II, une décision doit être justifiable par des faits,
        pas par un indice.
\end{itemize}

\subsection{Le tableau de bord climatique}

Copernicus et Open-Meteo alimentaient un panneau de variables climatiques. Le
problème n'était pas la donnée, mais sa forme : les valeurs exposées étaient un
\emph{indice d'exposition de 0 à 100}, dérivé, et non des variables
climatiques en unités réelles (jours de canicule par an, degrés Celsius,
mètres par seconde). Un indice 0--100 n'est pas une donnée climatique : c'est
un jugement déguisé. Il tombait sous le coup de la règle sur les scores.

\subsection{La décoration 3D}

Les extrusions 3D des bâtiments n'apportaient aucune information : elles
montraient une forme sans dire si le bâtiment était exposé. C'était un coût de
performance et de complexité pour zéro contenu.

\subsection{La cause commune}

Ces quatre symptômes ont une racine unique : \textbf{le produit n'était pas
organisé autour de la seule chose qu'il faisait et que personne d'autre ne
faisait} — la jointure entre exposition et vulnérabilité, à l'échelle du
bâtiment. L'effort était réparti sur quatre connecteurs, une couche de
notation, un tableau de bord et de la 3D, plutôt que concentré sur cette
jointure.

\section{La décision de recentrage}

Le 25 août 2026, une session de cadrage a produit une réorientation complète,
documentée dans \fichier{docs/workflow/concept.md} et
\fichier{01-brainstorm-notes.md}. Elle tient en cinq points.

\begin{enumerate}[leftmargin=1.6em]
  \item \textbf{Le livrable est un contrat de données.} L'interface devient un
        outil de vente et de démonstration, pas le produit.
  \item \textbf{Le score composite est mort.} Faits et provenance uniquement.
  \item \textbf{Le périmètre de la v1 est la jointure Géorisques $\times$
        BDNB.} Copernicus et Open-Meteo n'ont pas d'usage assureur en v1.
  \item \textbf{Le champ de résolution devient le signal de qualité.} Puisque
        le score disparaît, c'est lui qui porte l'intégrité du contrat.
  \item \textbf{Le code mort est supprimé, pas mis en sommeil.}
\end{enumerate}

\section{L'exécution : supprimer}

La conséquence la plus visible du recentrage est une phase de suppression.

\begin{table}[H]
\centering
\small
\caption{Suppressions de la phase de recentrage, par catégorie.}
\label{tab:suppressions}
\begin{tabularx}{\textwidth}{@{}l X@{}}
\toprule
\textbf{Catégorie} & \textbf{Éléments supprimés} \\
\midrule
Connecteurs backend &
  \fichier{copernicus.py}, \fichier{era5\_land.py}, \fichier{flood.py},
  \fichier{fire\_danger.py}, \fichier{heat\_stress.py}, \fichier{wind\_risk.py},
  \fichier{seasonal.py} \\
\addlinespace
Routes backend &
  \fichier{climate.py}, routes \api{/diagnostic/copernicus/*} \\
\addlinespace
Schémas &
  \fichier{zone\_assessment.py}, schéma \code{RecommandationsIA} et le champ
  \code{recommandations} \\
\addlinespace
Scoring &
  \fichier{risk\_model.py} et toutes ses références, \fichier{app/scoring/} \\
\addlinespace
Dépendances &
  \code{cdsapi}, \code{xarray}, \code{netCDF4} retirés de
  \fichier{requirements.txt} \\
\addlinespace
Tests &
  \fichier{test\_climate\_dashboard.py}, \fichier{test\_copernicus\_*.py} (4),
  \fichier{test\_trajectoire\_brute.py} \\
\addlinespace
Composants frontend &
  \fichier{CopernicusPanel.tsx}, \fichier{CopernicusStatusBanner.tsx},
  répertoire \fichier{jumeau/} (moteur de scène, LIDAR, recommandations),
  \fichier{Artisans.tsx} \\
\addlinespace
Styles &
  \fichier{copernicus.css}, \fichier{jumeau.css}, \fichier{artisans.css} \\
\bottomrule
\end{tabularx}
\end{table}

Deux règles de méthode ont été respectées pendant ces suppressions. D'abord,
elles constituent un \emph{change set dédié} : on ne supprime pas du code au
milieu d'une tâche de fonctionnalité, car cela rend la revue impossible. Ensuite,
le critère n'était pas \og c'est ancien \fg{} mais \og c'est inatteignable \fg :
avant de supprimer, on vérifie qu'aucune route, aucun import, aucune référence
ne pointe vers l'élément. Un code partagé et vivant, comme le module de
recommandations du répertoire \fichier{jumeau/}, a été \emph{déplacé} plutôt que
supprimé.

\section{Le compromis assumé : ce qui n'a pas été supprimé}

Deux éléments ont été conservés malgré la logique de recentrage, et il faut le
dire.

\begin{description}[leftmargin=0em,style=nextline]
  \item[La carte.] Elle n'est pas de la décoration : le souscripteur fait du
  jugement spatial. Mais elle a été ré-orientée : au lieu d'afficher les
  couches sources, elle affiche la jointure. L'empreinte du bâtiment est
  colorée par le résultat du test géométrique, accompagnée d'une fiche de
  vulnérabilité.
  \item[Les projections climatiques.] La direction produit initiale les
  excluait de la v1, mais la non-stationnarité (chapitre~\ref{ch:metier}) est
  l'argument de fond du produit. Les projections sont donc revenues, sous une
  forme différente : non plus comme un tableau de bord de variables, mais comme
  une extension honnête du contrat, avec de vraies unités et une étiquette de
  résolution.
\end{description}

\section{Enseignement}

La leçon de ce chapitre est méthodologique. Un projet technique n'est pas
\og bien avancé \fg{} parce qu'il contient beaucoup de code. Il est bien avancé
quand il devient difficile de lui retirer quelque chose sans le casser. Le
recentrage a réduit le périmètre fonctionnel de manière significative et
augmenté la valeur du livrable — parce que la valeur n'était pas dans le
nombre de fonctionnalités, mais dans l'exactitude de la jointure et dans la
traçabilité de chaque fait.

\begin{figure}[H]
\centering
\begin{tikzpicture}[
  scale=0.95, transform shape,
  box/.style={draw,rounded corners=2pt,align=center,inner sep=5pt,font=\small},
  removed/.style={box,draw=typhoonRed,fill=typhoonRed!6,text width=2.6cm},
  kept/.style={box,draw=typhoonGreen,fill=typhoonGreen!8,text width=2.6cm},
  core/.style={box,draw=typhoonBlue,very thick,fill=typhoonBlue!10,text width=3.4cm}
]
\node[removed] (c1) at (0,2.4) {Connecteurs climatiques};
\node[removed] (c2) at (0,1.2) {Couche de notation};
\node[removed] (c3) at (0,0) {Tableau de bord climat};
\node[removed] (c4) at (0,-1.2) {Extrusions 3D décoratives};
\node[removed] (c5) at (0,-2.4) {Rapport narratif IA};

\node[core] (core) at (5.6,0) {\textbf{La jointure}\\[3pt]
  Aléa Géorisques $\times$ Vulnérabilité BDNB\\[3pt]
  {\footnotesize par bâtiment, avec provenance}};

\node[kept] (k1) at (11.2,1.8) {Carte de la jointure};
\node[kept] (k2) at (11.2,0.6) {13 aléas normalisés};
\node[kept] (k3) at (11.2,-0.6) {Contrat JSON canonique};
\node[kept] (k4) at (11.2,-1.8) {Moteur de scénarios TRI};

\foreach \x in {c1,c2,c3,c4,c5} \draw[-{Latex[length=2mm]},typhoonRed,dashed] (\x) -- (core);
\foreach \x in {k1,k2,k3,k4} \draw[-{Latex[length=2mm]},typhoonGreen] (core) -- (\x);

\node[font=\small\itshape,color=typhoonRed] at (2.6,3.3) {retiré du périmètre};
\node[font=\small\itshape,color=typhoonGreen] at (11.2,3.3) {conservé / construit};
\end{tikzpicture}
\caption{Le recentrage en une image. Tout ce qui ne sert pas la jointure
aléa $\times$ vulnérabilité est retiré du périmètre ou relégué à la couche de
présentation.}
\label{fig:recentrage}
\end{figure}
